\section{Introduction}

The gelation of polymer solutions and colloidal suspensions — the transition from a viscous fluid to a viscoelastic solid — is characterised by a well-defined gel-point $\pcprime$ at which the loss tangent $\tan\delta = G''/G'$ becomes frequency-independent and equal to unity \citep{Chambon1987,Winter1987}. In the Winter--Chambon framework this corresponds to a phase angle $\delta = 45°$ and a power-law mean squared displacement (MSD) $\langle r^2(t) \rangle \propto t^{0.5}$, measurable by passive microrheology through the Generalised Stokes--Einstein Relation (GSER) \citep{Mason1995,Mason2000,Squires2010}. The gel-point is routinely used as the defining characteristic of a material system, and a considerable body of literature has established its dependence on polymer concentration, cross-link density, and temperature \citep{Rubinstein2003,Larson1999}.

What is rarely asked, however, is whether the gel-point is also a function of \emph{how the network forms} — that is, of the spatial growth dynamics of the cluster itself. In colloidal gels, crystallising proteins, and biopolymer networks, cluster growth is not random; it is governed by nucleation and propagation rules that produce morphologies ranging from fractal aggregates to compact, space-filling domains \citep{Trappe2001,Gibaud2012,Ng2016}. Whether these different morphologies shift the gel-point, and whether that shift can be controlled by a single parameter, has not been systematically established.

In this paper we use three-dimensional random walk (RW) simulations on site-percolation lattices to address this question directly. The approach is motivated by passive microrheology: probe particles in a gelling medium diffuse through the \emph{void} (fluid) phase of the network, not through the network itself \citep{deBruyn2013,Rich2011,Oppong2006}. We therefore place walkers on the unoccupied sublattice of an $L \times L \times L$ cubic lattice, sweeping the occupation probability $p$ from 0 to 1 and tracking the anomalous diffusion exponent $\alpha(p)$. The gel-point $\pcprime$ is identified by $\alpha = 0.5$, which corresponds to $\delta = 45°$ via the GSER. This places the model firmly within the Winter--Chambon framework for critical gelation \citep{Winter1987,Chambon1987,Muthukumar1989}.

We present three interconnected results. First, for randomly generated percolation networks (Bernoulli site percolation), the gel-point coincides quantitatively with the loss of spanning connectivity in the void sublattice: $\pcprime \approx 1 - \pc \approx 0.683$, where $\pc = 0.3116$ is the 3D site percolation threshold. This is not a coincidence — it is a direct consequence of the walkers probing the void-phase topology, and it provides a quantitative link between the rheological gel-point and a well-characterised geometric phase transition. Second, replacing random growth with adjacency-constrained templated growth — where new occupied sites must be adjacent to existing cluster members, mimicking nucleation-and-growth kinetics — shifts the gel-point by $\Delta\pcprime = +0.20$ to $\pcprime \approx 0.886$, and qualitatively changes the nature of the transition. Third, a single nucleation-density parameter $\kappa$ interpolates continuously between these two limits, providing monotonic control over $\pcprime$ across a range $\Delta\pcprime \approx 0.31$. We complement these structural results with GSER-derived frequency-domain spectra $G'(\omega)$ and $G''(\omega)$, confirming the $\delta(\omega) = 45°$ signature at $\pcprime$ for both gelation classes.

Together these results demonstrate that the gel-point of a percolating network is not fixed by density alone — it is set by the growth dynamics of the network-forming phase, and is tunable via the nucleation density of the growing cluster.

